\scanpage{7}
\mspage{6}\hfill\mspage{3}
for every diagram
\diagram{\[
\left\{\begin{gathered}\begin{tikzcd}[row sep=large]
\deleted{A}\inserted{X}&\\
{|X|}\arrow[r,"{\substack{\alpha\\\sim}}"]&E
\end{tikzcd}\end{gathered}\right.
\]}
with $X\in\mathrm{Ob}\,\mathrm{Esp}_{\mathcal M_*}$, $\exists$ a unique \deleted{$Y$ and} arrow $f:X\mathbin{-}Y$ in $\mathrm{Esp}_{\mathcal M_*}$, which lifts $\alpha$ (``Transport of structure'').

\uline{Proposition} In the category $\mathrm{Esp}_{\mathcal M_*}$ of $\mathcal M_*$-tame spaces, there are
\deleted{fiber products \inserted{\deleted{(of two objects} over a third)}, and products (of two objects), and the functor $X\mathrel{\mapstochar\relbar}|X|$ commutes with them. If one assumes
\[
(M5)\qquad R_0\ne\varnothing
\]
then there \deleted{functor} is also a terminal object (which has one point) and the functor}
\inserted{finite limits}
and the functor $X\mathrel{\mapstochar\relbar}|X|$ commutes with them.
\margin{\unsure{triviality of} \unsure{followed} by a \illegible{1} \unsure{and of the} \illegible{2} of the \unsure{functors}}

\uline{Proposition} If for \uline{one} $n\in\mathbb N^*$ one has \deleted{$\unsure{R}$} $R^n\in\mathcal M_n$, then this is true for $\forall n$, and even for $\forall I$ finite \inserted{in}
\[
R^I\in\mathcal M_I.
\]
Let us consider the axiom (\uline{optional})
\[
\deleted{\unsure{M5}}(M0)\qquad R\in\mathcal M_1\qquad(\Longrightarrow\forall I\in(\mathrm{Ens\ finite}),\ R^I\in\mathcal M_I)
\]
\margin{When one does not have $(M0)$, but assumes $(M5)$, then the union of the $A\in\mathcal M_1$, means that $\unsure{R'}$ is \unsure{obtained} \unsure{as} \illegible{2} \unsure{of} $\unsure{R}$ \unsure{filtering} \unsure{union} \illegible{2} $\unsure{\mathcal E|\Sigma}$.\par
\unsure{One has} \unsure{these} \unsure{stable} \unsure{sets} \illegible{2} the axiom $(M1)$ \unsure{is expressed} \illegible{1} \unsure{the influence} \unsure{of} \unsure{functors} \illegible{2}.\par
\unsure{as} $(M5)$ the expression \unsure{space} \unsure{equipped with} \unsure{the} $\unsure{R}$ \unsure{also} \unsure{empty} \unsure{annulled}.}
\margin{\unsure{object} \unsure{abs.} \unsure{Cat}\uline{\unsure{f}}}
